\section{Lezione 19} \hfill 21/04/2026\\
Le piantine ordinate arrivano in cantiere: operativamente non è possibile piantarle in uno stesso giorno. Le rese quotidiane ad operaio sono solitamente di 80-100 piante se a radice nuda e 200-300 se in pane di terra.\\
Il materiale dev'essere quindi conservato:
\begin{itemize}[nosep]
    \item a pane di terra: dev'esserci un luogo all'ombra, e dotato di un impianto di irrigazione (in vaso tendono a seccarsi velocemente);
    \item a radice nuda: le piante arrivano a mazzi di 25-50 piantine (avvolte in sacchi di iuta), e le si pongono in taiola (si crea un solco dove vengono messi i mazzi, e vengono poi ricoperti di terra). 
\end{itemize}
Le piantine devono essere portate nel luogo dove viene fatto il rimboschimento:
\begin{itemize}[nosep]
    \item quelle a pane di terra vengono messe in cassette e trasportate;
    \item quelle a radice nuda vengono imbozzimate (le si pongono in un bugliolo con una fanghiglia di terra ed acqua). La poltiglia permette sia di idratare l'apparato radicale e sia di proteggere le radici dalla luce solare.
\end{itemize}
Per quanto riguarda le piantine a radice nuda, occorre ridurre al minimo lo shock da trapianto. Per fare ciò, la regola dice che:
\begin{enumerate}[nosep]
    \item occorre scavare una buca da 40x40x40 cm di lato, anche se la piantina è piccola. Il cotico erboso, la terra scavata e le eventuali pietre presenti vengono messe da parte;
    \item dopo aver scavato la buca, si forma una piccola montagnola all'interno di essa;
    \item la piantina viene poi presa e posta al di sopra della montagnola;
    \item successivamente si riempe la buca con terra, senza compattarla;
    \item la piantina viene poi presa e tirata con cura su, fino al livello desiderato: in questo modo le radici vengono messe in posizione corretta verso il basso;
    \item dopo aver scelto la posizione corretta, si compatta bene il suolo: le radici non sopportano l'aria;
    \item la lettiera viene posta dove si trovava, con le radici all'aria (in modo che muoia). Le pietre invece si pongono intorno (soprattutto se il sito è freddo), in modo da scaldare la pianta.
\end{enumerate} 
Sarebbe meglio scavare la buca prima (rispetto che farlo al momento della semina), portando così alla luce/secco/aria eventuali patogeni: in questo modo viene svolta una specie di sterilizzazione (es. per piantare in primavera, si scava la buca in autunno).\\
C'è il rischio che la buca venga coperta nuovamente (es. da erosione, neve): dal punto di vista operativo è quindi più conveniente scavare subito prima.\\
Il rischio pratico più frequente durante un rimboschimento è la presenza di nidi di vespe a terra durante lo scavo.\\
Il colletto della piantina viene posto a diversi livelli rispetto al suolo, in relazione alle caratteristiche di umidità stazionali:
\begin{itemize}[nosep]
    \item se c'è poca acqua, conviene fare una conca, in modo da captare la rugiada mattutina;
    \item se le condizioni sono normali, si pone a livello del suolo;
    \item se c'è eccesso d'acqua, si crea una collinetta sopra il livello del suolo.
\end{itemize} 
Questa operazione viene fatta completamente a mano, soprattutto per l'impossibilità ad accedere con mezzi meccanici.\\
Per quanto riguarda il postime in pane di terra, la buca può essere appena più grande del vaso. Conviene comunque togliere lo strato erbaceo sopra per almeno 40x40 cm, in modo per evitare almeno all'inizio la competizione delle erbe con la piantina.
\subsection{Post-rimboschimento}
\subsubsection{Protezione dal brucamento}
Dopo aver fatto il rimboschimento, occorre proteggerlo dal brucamento.\\
Le piantine del vivaio sono particolarmente apprezzate dai selvatici: sono ricche di sostanze nutritive e per questo più appetibili.\\
Le protezioni possono essere:
\begin{itemize}[nosep]
    \item singole alle piantine, fatta mediante i shelter, ovvero tubi (cilindrici, quadrati o rettangolari) che avvolgono la piantina. Devono essere alti più dell'altezza di brucamento del selvatico (considerando anche lo spessore della neve). Gli shelter tendono a coricarsi, necessitando quindi di tutori (soprattutto dove c'è molta neve, che si tende ad evitarne l'utilizzo). Possono essere:
        \begin{enumerate}[nosep]
        \item chiusi: proteggono completamente la piante, ``spingendo'' la pianta verso l'alto. Ci può essere il problema delle elevate temperature in estate; inoltre, possono diventare luoghi adatti per le vespe per la creazione di nidi. 
        \item aperti: come fossero una specie di rete. La pianta tende a far crescere i rami dentro le maglie delle reti. 
        \end{enumerate}
    Per entrambi i casi, occorre rimuovere lo shelter e smaltirlo regolarmente.
    \item protezioni comuni (recinzioni): sono molto efficienti, ma costano molto (prima del COVID costavano 20 euro a metro lineare).\\
    Le recinzioni devono:
    \begin{enumerate}[nosep]
        \item resistere fino a quando le piante non escono dal rischio di brucamento (10-20 anni);
        \item essere alte almeno 1.5-2 metri, con delle maglie che non permettono il passaggio dei selvatici;
        \item prevedere un ingresso;
        \item se la recinzione non è troppo grossa, l'animale la vede e ci gira attorno;
        \item non essere rettilinee, per non portare il selvatico a saltarci dentro;
        \item non avere forme concave, per lo stesso motivo;
        \item essere inferiori all'ettaro, perché il perimetro aumenta notevolmente
    \end{enumerate}
\end{itemize}
In definitiva, considerando i pregi e difetti, è ancora meglio avere più brucamento rispetto all'utilizzo dello shelter.\\
Un'altra forma di protezione sono i repellenti, che vengono spennellati o spruzzati sulla piantina: in questo modo l'appetibilità dei selvatici si riduce. È molto usato in Austria, dove ci sono forti conflitti tra i cacciatori ed i forestali.
\subsubsection{Competizione erbacea}
Le piantine devono essere protette anche dalla competizione con le specie erbacee, soprattutto in pianura.\\
È possibile fare una pacciamatura andante con un disco pacciamante, di almeno 40 cm di diametro. È poco efficiente se le erbe adiacenti crescono molto (ed infatti non viene svolto in bosco). Solitamente sono di materiale plastico (da smaltire) o di materiale biocompostabile.\\
La pacciamatura (in generale) non conviene, è forse meglio andare a ripulire.\\
Al primo anno dal rimboschimento, dove le piantine sono fragili, conviene fare delle ripuliture (soprattutto se non ci sono pacciamature). Vengono svolte con decespugliatori, da evitare quelli a filo, ma sono da preferire quelli a disco. Il raggio di azione attorno alla piantina o al gruppo di piantine dev'essere almeno l'altezza delle piante erbacee attorno a loro.\\\\
Alla fine del primo anno del rimboschimento, ogni pianta ha un ruolo: ci si aspetta che ogni pianta sopravviva. Purtroppo, non è così e molte piante muoiono per diversi motivi. La quantità di piante tollerate (che possono morire) nelle condizioni migliori sono inferiori al 3\%, mentre nelle peggiori pari al 10\%. Se il numero è maggiore, occorre risarcire. Solitamente il collaudo del rimboschimento viene fatto al terzo anno: in quel momento, le piante morte devono essere inferiori alle tolleranze previste dal contratto.\\
I risarcimenti vengono fatto con le stesse specie usate per il rimboschimento. Le piante possono essere prese da un avanzo del rimboschimento: se le piante erano in pane di terra occorre verificare che le radici non siano spiralate e se erano a radice nuda occorre andare a vedere che siano sopravvissute nella tagliola. Eventualmente occorre andare ad acquistare nuovamente in vivaio (anche in questo caso si trattano di riserve).\\
Se dopo 3 anni ci sono più fallanze del progetto iniziale previsto, allora sono stati fatti degli errori durante la progettazione, realizzazione e/o manutenzione.\\
Il rimboschimento comincia a crescere, e le piante (soprattutto in un sesto regolare) hanno una certa distanza tra di loro: la crescita avviene in modo costante e simile (non ci sono piante che ``spingono'' altre verso l'alto). In questo caso tutte le piante sono forti e dominanti (non hanno competitori intorno): questo fino a che non si toccano.\\
La competizione laterale, prima di cedere, le porta a crescere più di quelle accanto: si generano così piante filate. In un rimboschimento, gli sfolli ed i diradamenti sono necessari, e non auspicati: occorre evitare che le piante si disequilibrino (si cerca di ottenere il rapporto di snellezza minore a 100 ed una chioma intorno al 50\% del fusto). Invece, in un popolamento naturale questi trattamenti possono essere evitati, in quanto è la natura a determinare quali alberi sono i vincitori della competizione.\\
Se non vengono svolti diradamenti, si ottengono piante filate, instabili ed in competizione (gli apparati radicali sono limitati).\\
La tempesta Vaia ha colpito soprattutto i rimboschimenti: se ci fossero state utilizzazioni corrette, il danno sarebbe stato minore.
\subsubsection{Potatura}
Una operazione disponibile, ma poco fatta, è la potatura da legno: l'obiettivo è la creazione di un toppo (di 2-6 metri a seconda della specie) esente da rami, in modo da ottenere legno di prima qualità.\\
La concezione della potatura da legno è completamente diversa dagli altri tipi di potatura:
\begin{itemize}[nosep]
    \item estetica (anche detta ``arte topiaria''): l'obiettivo è quello di avere una bella chioma;
    \item produttiva:
    \begin{enumerate}[nosep]
        \item per frutti: inducendo lo stress alla pianta, la si porta a rinnovarsi. Si vuole ottenere una produzione massima, costante e di alta qualità del frutto, oltre che semplificare il processo di raccolta;
        \item per legno: non si vuole ottenere nodi al toppo.
    \end{enumerate}
\end{itemize}
Potare un albero occorre scegliere un criterio adatto per ogni singola pianta (le architetture sono diverse).\\
La potatura da legno ha delle regole fondamentali:
\begin{itemize}[nosep]
    \item non potare mai nei primi 2-3 anni dopo l'impianto: la pianta deve fare radici e diventare forte; 
    \item non potare mai più del 50\% della chioma (nelle giovani piante il cimale è una chioma e non è fusto). Nel caso ci fosse una eccessiva potatura, la pianta reagisce facendo uscire nuovi rami;
    \item potare sempre un fusto che sia più piccolo di 10 cm di diametro (non DBH, ma dove si pota): tendenzialmente in quei 10 cm non legno di interesse d'utilizzazione (es. midollo, legno giovanile);
    \item non tagliare mai un ramo che sia più grosso di 3 cm: se piccola, la ferita chiude subito, solitamente in una stagione vegetativa;
    \item quando si taglia un ramo di primo ordine (inserito nel fusto), occorre rispettare il cercine (anello creato attorno al punto di inserzione del ramo): la pianta ha maggiori difficoltà a ricoprirlo. 
\end{itemize}
La potatura può essere fatta tutto l'anno. È sconsigliabile in primavera, sopratutto a stagione vegetativa appena iniziata: la pianta è molto reattiva. Il momento ideale è verso la fine della stagione vegetativa o nell'inverno prima del gelo.\\
Il modello di potatura per le conifere consiste nel rimuovere i palchi dal basso, fino ad ottenere un toppo utile. 